\ContentDes{（二）研究内容：}
% intentional line break

\NsfcNote{（提纲不做限制，请按照研究工作的自身逻辑撰写。应提炼出特色与创新点、年度研究计划）}

\section{拟解决的关键科学问题}

面对未来网络需求的形式和挑战，结合国内外研究的现状，
本项目在开发基于雾计算和移动点对点的流媒体传输及播放技术中，拟解决如下三个相互关联的关键问题：

\begin{enumerate}
\item \emph{如何开发适用于各种雾计算设备的软件及设备间的通信协议？}
不同于现有的台式计算机和服务器，雾计算设备往往只有有限的计算和存储性能。
现有的适用于服务器和个人电脑的流媒体应用程序往往很难在雾计算设备（如图2，例如路由器、机顶盒）得到应用。
我们必须开发低负载的、适合于雾计算设备的软件。同时，雾计算设备往往会采用各异的硬件架构（例如X86，arm等）
和不同的操作系统（例如linux，android等）。如何进行跨平台开发，使得软件能够兼容上述的雾计算设备，
将是不得不解决的问题。

\item \emph{如何开发适用于各种智能手机的移动点对点流媒体应用？}
雾计算依然需要部署硬件设备。在某些实际场景（例如在地铁车厢、公共汽车上），
部署雾计算设备就显得很困难。这时可以利用智能手机自身的资源建立Ad hoc的Wi-Fi网络提供流媒体服务。
例如，一台智能手机可以通过3G网络下载到流媒体视频，然后通过Wi-Fi分享给临近的其它智能手机。
然而，智能手机的性能有限，而且型号繁多，种类参差不齐。
如何保证智能手机在分享视频的同时自身性能不受影响，依然能够提供优秀的用户体验，也是必须考虑的挑战。

\item \emph{如何构建流媒体网络的拓扑结构？}
流媒体的拓扑结构是一直一个重要的研究方向。好的流媒体拓扑结构算法要满足如下条件。
首先，随着视频资源的增多，算法要考虑到多视频源、多频道的问题。
直接套用传统的算法很有可能在某处网络资源瓶颈造成拥堵，从而破坏了用户体验。
其次，好的算法也要保证在现有的资源下最大限度的提升用户体验，尽可能的缩短端到端的延迟。
最后，好的算法还需要考虑到网络节点的变化。
移动设备和雾计算设备可能会不时地上线和离线，算法必须能够应对这种变化，
在网络资源改变的时候重新调整网络拓扑结构，继续稳定地向用户提供服务。
\end{enumerate}

\section{研究内容}

结合以上分析，本项目在开发基于雾计算和移动点对点的流媒体传输及播放技术中，将主要研发如下内容：
\begin{enumerate}
\item \emph{在网络中大规模部署“雾计算”设备研究。}
我们将开发安装在雾计算装置上的软件，使各种装置可以为流媒体的传输贡献资源。
对于不同的设备，我们的软件应该能做到跨平台、跨操作系统的支持。
我们将搭建实验平台来测试雾计算装置的可靠性。
最终我们将实现商业化服务，用户可以将我们的软件安装在自己的雾端设备上
（或者直接购买我们的雾计算设备）来享受流媒体服务。
对于内容提供商来说，雾端设备能大大降低提供服务的成本，
只要少量的网络资源（云端服务器或者自己拥有服务器）就能服务大量的用户。
另外，在网络中大规模部署“雾计算”设备时，不仅使得网络能有效传输流媒体，
在某些设备失效的情况下，网络依然能有效工作。

\item \emph{在实际应用场景下“雾计算”和移动点对点流媒体技术融合研究。}
我们也将开发安装在智能手机上的应用，使得智能手机可以作为一个网络热点向周围的人分享流媒体。
手机应用将保证在不影响手机的用户体验下，向周围的用户提供服务。
智能手机通过网络下载到流媒体视频，利用手机自身的资源建立Ad hoc的Wi-Fi网络，
帮助其它手机用户。我们将研究在实际应用场景下将“雾计算”和移动点对点流媒体技术相结合，
在现实场景（例如地铁，机场，商场等人流密集区）验证技术的可靠性和成熟性，并最终实用化。

\item \emph{流媒体数据传输的优化路由算法研究。}
我们会开发优化流媒体数据传输的路由算法，以最大限度地合理利用网络资源、减小延迟、提升用户体验。
该算法将统筹考虑多视频源、多频道的情况，避免在网络某一瓶颈形成拥堵。
同时该算法将优先考虑用户体验，最小化延迟。
优化算法也将适应雾计算设备网络的动态变化，
在某些设备下线或失效的情况下能及时调整网络结构，继续稳定地提供服务。
\end{enumerate}

\section{拟采取的研究方案及可行性分析}

\subsection{拟采取的技术路线}

\begin{enumerate}
\item 雾计算设备的软件开发：我们将开发跨平台、跨操作系统的软件来实现利用雾计算设备来支持流媒体的传输。
\item 雾计算设备的实验平台：我们将用不同种类的雾计算设备（传统的路由器及新的智能路由器）
首先搭建固定拓扑结构的网络，验证雾计算设备软件的可靠性。
\item 移动点对点流媒体的软件开发：我们将开发利用手机自身的资源建立Ad hoc的Wi-Fi网络的软件，
在不影响手机的用户体验的前提下，向其它智能手机传输流媒体视频。
\item 流媒体数据传输的优化路由算法研究：我们将通过数学建模来捕捉网络的关键信息，
然后设计出能够在多视频源、多频道的条件下，能确保用户体验的优化路由算法。
我们将给出理论的证明，确保算法有好的性能。
\item 优化算法的仿真：我们将使用仿真软件来验证优化算法的有效性，研究在不同网络条件下网络性能的变化。
\item 综合的流媒体服务实验平台：在实验室搭建全要素的实验平台。平台将包括视频源服务器，雾计算设备和移动端设备。
我们将通过实验来验证仿真结果，并为最后的商业化产品打下基础。
\item 实际应用场景下的实验：我们将完善的雾计算解决方案及智能手机的应用，并推出测试版的服务，
用户可以在香港试用我们的产品。在通过实际应用场景下的实验之后，我们的产品将最终商业化并向用户提供服务。
\end{enumerate}

\subsection{可行性分析}

\begin{enumerate}
\item \emph{研究内容具体明确，研究方案切实可行。}
本项目主要关注应用雾计算技术和移动点对点技术提供流媒体服务。
在理论方面，我们已经有了支持传统媒体点播（Video-on-demand）、流媒体（Live Streaming）
及点对点文件分享（P2P）算法设计以及相关验证过程作为参考。
在实践方面，我们之前研究的网络电视和基于无线网络的室内定位技术已经实现商业化部署，
我们非常熟悉相关系统的开发。我们在以往研究中搭建的测试平台为我们今后的实验打下了良好的基础。

\item \emph{项目组具有承担本研究的能力。}
陈双幸教授领导的多媒体实验室（MTrec：Multimedia Technology Research Center）
自创建以来一直致力于无线网络和多媒体网络技术的研究，在相应的研究方向上都有一定程度的积累。
我们在相关领域上的期刊和会议论文超过150篇，绝大多数都是ACM/IEEE的会议及期刊。
我们还承担了香港政府以及工业界的多个大型研究项目，相关的研究成果已经转化为商业化应用，
开始服务香港市民。
\end{enumerate}

综上所述，我们已经拥有无线网络和多媒体服务研究的相关经验，为本项目的研究打下了良好的基础。
课题组的主要成员有着丰富的研究和开发经验，并且在以往已经取得了良好的成果。
通过详细的模型构建和算法设计，结合仿真及实验平台，我们可以有效的测试、验证和改善解决方案，
确保在实际应用场景中成功提供服务。
最终，我们完全能够完成既定的研究目标，在发表相关论文的同时部署商业化的解决方案。

\section{年度研究计划及预期研究结果}

\subsection{年度研究计划}
\begin{itemize}
\item 第一阶段（半年：2018年1月1日至2018年6月30日）：
文献综述，深入研究雾计算和移动点对点的传输协议和设备的驱动程序/操作系统。
针对多媒体传输网络的特性，提出项目总体研究计划方案。完成流媒体传输服务理论模型构建。\\
资金使用：劳务费30万。

\item 第二阶段（半年：2018年7月1日至2018年12月31日）：提出多媒体传输网络的构造算法，
并完成理论证明和仿真验证。并从用户体验、端到端延迟和稳定性等方面分析和验证对于多媒体服务的影响。\\
学术目标：完成流媒体数据传输的优化路由算法研究。撰写论文并投稿会议或期刊。
资金使用：劳务费30万，材料费15万用于购买计算机软件。

\item 第三阶段（半年：2019年1月1日至2019年6月30日）：开发各个设备（服务器、雾端和移动端）的软件，
构建实验平台，完成现实条件下的系统测试和验证。开发商业产品的原型。 \\
学术目标：完成在实际应用场景下“雾计算”和移动点对点流媒体技术融合的研究，
完成在网络中大规模部署“雾计算”设备研究。开始申报相关专利。\\
资金使用：劳务费30万，材料费17万及设备费15万用于构建实验平台。

\item 第四阶段（半年：2019年7月1日至2019年12月31日）：总结研究成果，完善产品，完成研究报告的编写，
准备项目结题验收。\\
学术目标：参加投稿论文的学术会议，完善已有的成果。\\
资金使用：劳务费30万，差旅费12万及会议费8万用于参加学术会议。
\end{itemize}

\subsection{预期研究成果}

面对未来流媒体技术所面临的挑战和实际需求，当前的主流多媒体传输技术将面临网络资源捉襟见肘的挑战。
我们将利用雾计算和移动点对点技术所带来的优势，开发优化流媒体数据传输的路由算法，
使其能够支持大量用户同时观看高清晰度的流媒体视频，同时能够在部分网络设备无法工作时，
保证剩余的设备依然能够支持服务。我们将开发安装在雾设备和智能手机上的软件，
使得市场上现有的设备可以为流媒体的传输贡献资源，将研究项目所产生的新技术落实到现有的用户。

对于研究中带来的新成果，我们将撰写论文投稿在国际权威会议及期刊上。
对于新的技术，我们也将会申请专利。在研究过程中，我们也会培养和锻炼相关的博士研究生。

具体的预期目标如下：

学术指标：相关技术的论文将会在国际权威会议及期刊上投稿。
\begin{enumerate}
\item 国际会议或期刊发表高水平论文10篇或以上。
\item 申请发明专利3项以上。
\item 培养博士研究生3名。
\end{enumerate}

技术指标：相关技术可以大规模、商业化的部署。通过该技术，我们可以向区域乃至全国的用户提供媒体服务。
具体技术指标如下：
\begin{enumerate}
\item 支持2000个以上的用户同时观看视频。最终产品可以大规模地、商业化地部署在多种现有和未来的网络装置
（包括智能电话、平板电脑、网络路由器）中。使各种装置可以为流媒体的传输贡献资源。
\item 支持FHD视频，即支持3Mbps码率的流媒体流畅传输。在移动端，我们将开发新的软件架构，使得移动智能设备
（例如智能手机）能通过Ad hoc的WiFi模式向其它移动用户提供流媒体，从而节省移动网络（例如3G和4G）的带宽。
在雾设备端，我们将改进现有设备的驱动和操作系统，使其能满足支持FHD视频的要求。
\item 具有一定的健壮性，当5\%以内的雾设备无法正常工作时，网络仍然能依靠剩余正常的设备维持网络运转。
我们会开发优化流媒体数据传输的路由算法，以最大限度地合理利用网络资源、减小延迟、提升用户体验。
路由算法在发现雾设备无法工作时，将重新构建网络的拓扑结构，使得剩余的设备依然能够向全部用户提供服务。
\end{enumerate}
